% TOP_PROBLEM: {"id": 3088, "title": "Partitioning the Projective Plane", "category_id": 6, "category": {"id": 6, "name": "geometry", "display_name": "Geometry", "description": "Euclidean and non-Euclidean geometry, geometric structures.", "slug": "geometry", "order_index": 6, "created_at": "2026-07-31T15:26:25.671Z"}, "status": "open", "problem_number": "OPG-56328", "set_id": 12}
% FIRST_POSTED: 2026-10-01
% ENTRY_KIND: statement_only
% UnsolvedMath ID 3088; source catalogue code OPG-56328
% !TeX program = lualatex
% Definition and sources only; this file is not an LLM solution attempt.
\documentclass[11pt,a4paper]{article}
\usepackage[margin=25mm]{geometry}
\usepackage{fontspec}
\setmainfont{lmroman10-regular.otf}[BoldFont=lmroman10-bold.otf,
 ItalicFont=lmroman10-italic.otf,BoldItalicFont=lmroman10-bolditalic.otf]
\usepackage{amsmath,amssymb,mathtools}
\usepackage{unicode-math}
\setmathfont{latinmodern-math.otf}
\AtBeginDocument{\let\setminus\smallsetminus}
\usepackage{xurl,hyperref}
\hypersetup{colorlinks=true,urlcolor=blue}
\setcounter{secnumdepth}{0}
\setlength{\parindent}{0pt}
\setlength{\parskip}{5pt}
\setlength{\emergencystretch}{3em}
\begin{document}
\section*{Partitioning the Projective Plane}\label{3088}
{\small UnsolvedMath ID 3088 · OPG-56328 · Geometry · OpenGarden}
\subsection{Definitions and mathematical statement}
Write $\mathbb{RP}^2$ for the set of one-dimensional linear subspaces of
$\mathbb R^3$. Its elements are unoriented lines through the origin.
For an orthonormal basis $e=(e_1,e_2,e_3)$ define
\[
 O_e=\{\mathbb R x:x\ne0,\quad
                x=t_1e_1+t_2e_2+t_3e_3,\quad t_1,t_2,t_3\ge0\}.
\]
A set of lines $S$ is octahedral when $S\subseteq O_e$ for some basis $e$.
This agrees with the source's geometric definition: every line meets
the closed face $\operatorname{conv}\{e_1,e_2,e_3\}$ and its opposite in
the regular octahedron with vertices $\pm e_i$. Zero coordinates are allowed,
so face edges and vertices are included.

Call $S$ weakly octahedral if each three-element subset of $S$ is
contained in some $O_e$, where the basis may depend on the triple.
The conjecture states that if
\[
 \mathbb{RP}^2=S_1\mathbin{\dot\cup}S_2\mathbin{\dot\cup}S_3
                         \mathbin{\dot\cup}S_4
\]
and every $S_i$ is weakly octahedral, then every $S_i$ is octahedral.
Thus for each part there should be one orthonormal basis working for
all its lines simultaneously. Different parts may use different bases.

The partition must cover every projective line, with disjoint parts;
no measurability or closedness assumption is imposed. The local triple
condition alone, for an arbitrary subset not occurring as a part of such
a partition, is not the assertion. Directions are identified up to all
nonzero scalar multiples, so the choice of orientation of representatives
cannot affect either condition.
\subsection{Short English statement}
If the real projective plane is partitioned into four sets whose every triple fits in an octahedral face pair, must each entire set fit in one?
\subsection{Sources}
\begin{itemize}
\item[S1] ULAM.AI and the UnsolvedMath Contributors, supplied problems.json, record 3088 (OPG-56328), OpenGarden collection. Catalogue identity and source transcription; CC BY 4.0.
\url{https://huggingface.co/datasets/ulamai/UnsolvedMath}.
\item[S2] Open Problem Garden, Partitioning the Projective Plane. Original problem and discussion; the mathematical formulation below is expanded from the supplied source record.
\url{http://www.openproblemgarden.org/op/partitioning_the_projective_plane}.
\end{itemize}
\end{document}
